% !TeX program = xelatex
\documentclass[UTF8,a4paper,11pt,fontset=none]{ctexart}
\usepackage{geometry,fontspec,amsmath,enumitem,tabularx,array,xcolor,hyperref}
\geometry{left=1.7cm,right=1.7cm,top=1.25cm,bottom=1.25cm}
\setmainfont{Times New Roman}
\setCJKmainfont[AutoFakeBold=2.2]{SimSun}
\newCJKfontfamily\heiti[AutoFakeBold=2.0]{SimHei}
\hypersetup{colorlinks=true,linkcolor=black}
\setlength{\parindent}{0pt}
\setlength{\parskip}{0.10em}
\setlist[enumerate]{leftmargin=1.65em,itemsep=0.16em,topsep=0.08em}
\newcommand{\mm}{\,\mathrm{mm}}
\newcommand{\N}{\,\mathrm{N}}
\newcommand{\g}{\,\mathrm{g}}
\newcommand{\MPa}{\,\mathrm{MPa}}
\newcommand{\q}[2]{\item \textbf{问：#1}\par\textbf{答：}#2}
\begin{document}
\begin{center}
{\heiti\bfseries\zihao{2}勒勒车模型答辩快速稿}\par
{\zihao{-4}基于 350 mm 正式提交计算书｜只背结论、关键数字和证据}
\end{center}

\fbox{\begin{minipage}{0.96\linewidth}
\textbf{开场：}本作品按赛题 $40$ 袋满载设计，采用两根 $350\mm$、$22\times22\mm$ 方形空心主梁和 6 根横向构件。第 1、6 根为三角截面杆，第 2～6 根端部外伸；主梁强度与挠度采用 4 个等效加载点计算。模型质量 $203\g$，车轮为十二辐式，外径 $70\mm$。
\end{minipage}}

\section*{一、必须背熟的数字}
\begin{tabularx}{\linewidth}{>{\bfseries}p{3.0cm} X}
赛题与装载 & 每袋约 $500\g$，允许 $20\sim40$ 袋；本方案取 $40$ 袋，不捆绑布袋，除车轮外不得接触赛道或障碍物。\\
主要尺寸 & 主梁长 $350\mm$，截面外边长 $22\mm$，壁厚 $1.2\mm$；横向构件长 $250\mm$；车轮外径 $70\mm$，轮辐有效长度 $22\mm$。\\
荷载 & 货物重 $196\N$，模型自重约 $1.99\N$，总荷载 $197.99\N$。\\
主梁强度 & $\gamma_d=1.5$ 时 $9.89\MPa$、安全系数 $3.03$；$\gamma_d=2.0$ 时 $13.19\MPa$、安全系数 $2.27$。\\
挠度 & 理想静载 $1.31\mm$；实测左右 $5.6/5.8\mm$，平均 $5.70\mm$；校准设计挠度为 $8.55\mm$（1.5）和 $11.40\mm$（2.0）。\\
连接与坡道 & 单端横杆剪力 $47.78\N$；四螺丝水平牵引力 $36.9\N$、竖向支承力 $148.49\N$；坡道需求牵引力 $19.74\N$。\\
\end{tabularx}

\section*{二、高频提问与直接回答}
\begin{enumerate}
\q{实物到底是几根横杆？为什么计算只用了四个点？}{实物是 6 根横向构件；第 1、6 根为三角杆，第 2～6 根端部外伸。主梁整体强度和挠度把 6 根实物构件合并成 4 个等效加载点，这是计算简化，不是实物数量。}
\q{为什么主梁按 $350\mm$ 计算？}{$350\mm$ 是最终两根纵向空心主梁的实际长度。计算采用全长两端简支的统一模型，不利用外伸段可能带来的有利卸载。}
\q{总荷载怎样传递？}{货物和模型自重合计 $197.99\N$，经 6 根横向构件传给两根主梁，再传到车轴固定座、车轴、轴承和车轮。取 $\gamma_d=1.5$ 后，每根主梁设计荷载为 $148.49\N$。}
\q{$148.49\N$ 是不是平均单轮反力？}{不是。对称工况下单个后轮设计反力为 $74.25\N$；$148.49\N$ 是后轴总设计反力，也作为单桥一侧车轮的极端包络反力。}
\q{主梁强度是否满足？}{基础动力工况应力为 $9.89\MPa$，安全系数 $30/9.89=3.03$；严重冲击取 $\gamma_d=2.0$，应力 $13.19\MPa$，安全系数 $2.27$，均不低于 2。实物壁厚 $1.2\mm$ 也大于所需最小壁厚 $1.03\mm$。}
\q{为什么实测挠度比理论大？}{理论 $1.31\mm$ 是理想梁的静载结果，实测平均为 $5.70\mm$。差异主要来自胶接滑移、横向构件柔度和制作误差，因此反算等效弹性模量约 $1.38\mathrm{GPa}$，再用实测值校准得到 $8.55\mm$ 和 $11.40\mm$。}
\q{MIDAS 云图能不能直接证明最大应力？}{不能把现有截图当作可追溯的定量证据。云图用于说明端部和连接区的相对高值分布；最大应力、安全系数和挠度以统一手算、实测和试验为准。}
\end{enumerate}

\section*{三、老师继续追问时}
\begin{enumerate}[start=8]
\q{车轮怎么验算？}{车轮为十二辐式，外径 $70\mm$，轮辐有效长度 $22\mm$。按单轮极端包络取一根下部轮辐承担 $148.49\N$；1 次 $36$ 袋、2 次 $40$ 袋满载试跑后无开裂、松动、偏斜或卡滞。}
\q{四颗螺丝为什么还要算竖向力？}{因为动力小车连接端同时作为竖向支座。四颗 M4 螺丝同时承担 $36.9\N$ 水平牵引力和 $148.49\N$ 竖向支承力；单颗合力 $38.25\N$，剪应力 $4.35\MPa$，孔边承压 $3.19\MPa$。}
\q{坡道能不能通过？}{$4^\circ$ 坡道所需牵引力为 $19.74\N$；电机轴端理论牵引力为 $36.9\N$，所以总传动效率至少需 $53.5\%$，等效附着系数至少需 $0.20$。两次 $40$ 袋满载试跑均通过。}
\q{弯道如何保证不翻？}{取支承宽度 $300\mm$、重心高度 $120\mm$、速度 $0.30\mathrm{m/s}$、半径 $0.30\mathrm{m}$，有 $a_y/g=0.031<1.25=(B/2)/h_c$，整体抗倾覆裕度充足；实际重点是防止布袋侧滑。}
\q{为什么选择 40 袋？}{赛题允许 $20\sim40$ 袋，40 袋有效货物质量最大。模型质量 $203\g$、货物质量 $20000\g$，载质比为 $98.52$，且满载试跑通过。}
\q{试验有哪些证据？}{挠度实测为 $5.6/5.8\mm$；40 袋装货时间最长 $174\mathrm{s}$，距 180 秒上限余 $6\mathrm{s}$；两次 40 袋满载行驶时间均为 $34\mathrm{s}$，无开裂、松动、掉袋或非车轮部位碰撞。}
\q{为什么选方案 C？}{方案 C 的主梁—车轴固定座传力路径最清楚，关键节点少，制作误差更容易控制；综合强度、重量、荷载效率和制作工艺后作为最终方案基础。}
\end{enumerate}

\section*{四、三句话不要说错}
\begin{enumerate}
  \item 不要说“实物只有 4 根横杆”；正确说法是“实物 6 根，计算采用 4 个等效加载点”。
  \item 不要说“$148.49\N$ 是平均单轮反力”；对称平均单轮是 $74.25\N$，$148.49\N$ 是后轴总反力或单轮极端包络。
  \item 不要说“实测挠度小于理论值”或“MIDAS 直接证明最大应力”；实测 $5.70\mm$ 大于理想理论值 $1.31\mm$，云图只作定性分布说明。
\end{enumerate}

\vfill
\begin{center}\small\textbf{答题顺序：先结论，再数字，最后给计算或试验依据。}\end{center}
\end{document}
